\documentclass[a4paper,dvipsnames,twoside,12pt]{report}

% ----------------------------------------------------------------------
% General information for the title page and the page headings
\newcommand{\Title}{Statistical Practice in Epidemiology\\ with
     \includegraphics[height=1em,keepaspectratio]{../adm/Rlogo}\\
      Computer exercises}
\newcommand{\Tit}{SPE practicals}
%\newcommand{\Version}{2}
\newcommand{\Dates}{14--20 June 2018}
\newcommand{\Where}{IARC Lyon}
\newcommand{\Homepage}{\url{http://bendixcarstensen.com/SPE} }
% This years draw:
% > sample( c("B","K","J","E","M") )
% [1] "J" "B" "E" "K" "M"
\newcommand{\Faculty}{\begin{tabular}{rl}
Janne Pitk\"{a}niemi
  & Finnish Cancer Registry, Helsinki, Finland\\
  & \texttt{janne.pitkaniemi@cancer.fi}\\[1ex]
Bendix Carstensen
  & Steno Diabetes Center Copenhagen, Gentofte, Denmark \\
  & \& Dept. of Biostatistics, University of Copenhagen, Denmark\\
  & \texttt{b@bxc.dk} \\
  & \url{http://BendixCarstensen.com}\\[1ex]
Esa L\"{a}\"{a}r\"{a}
  & Department of Mathematical Sciences, University of Oulu, Finland\\
  & \texttt{Esa.Laara@oulu.fi} \\
  & \url{http://www.oulu.fi/university/researcher/esa-laara}\\[1ex]
Krista Fischer
  & Estonian Genome Center, University of Tartu, Estonia. \\
  & \texttt{Krista.Fischer@ut.ee} \\[1ex]
Martyn Plummer
  & International Agency for Research on Cancer, Lyon, France\\
  & \texttt{plummerm@iarc.fr} \\[1ex]
\end{tabular}}

\setcounter{tocdepth}{1}
% Stuff to put in at the begiing of the document
%----------------------------------------------------------------------
% Packages
\usepackage[utf8]{inputenc}
\usepackage[T1]{fontenc}
\usepackage[english]{babel}
\usepackage[font=it,labelfont=normalfont]{caption}
\usepackage[colorlinks,urlcolor=blue,linkcolor=red,citecolor=Maroon]{hyperref}
\usepackage[ae,hyper]{Rd}
\usepackage[dvipsnames]{xcolor}
\usepackage[super]{nth}
\usepackage[noae]{Sweave}
\usepackage{makeidx,floatflt,amsmath,amsfonts,amsbsy,enumitem,dcolumn,needspace}
\usepackage{ifthen,calc,eso-pic,everyshi}
\usepackage{booktabs,longtable,rotating,graphicx,subfig}
\usepackage{pdfpages,verbatim,fancyhdr,datetime,afterpage}
\usepackage[abspath]{currfile}
\renewcommand{\textfraction}{0.0}
\renewcommand{\topfraction}{1.0}
\renewcommand{\bottomfraction}{1.0}
\renewcommand{\floatpagefraction}{0.9}
\definecolor{blaa}{RGB}{99,99,255}
\DeclareGraphicsExtensions{.png,.pdf,.jpg}
% Make the Sweave output nicer
\DefineVerbatimEnvironment{Sinput}{Verbatim}{fontsize=\small,fontshape=sl,formatcom=\color{BlueViolet}}
\DefineVerbatimEnvironment{Soutput}{Verbatim}{fontsize=\small,formatcom=\color{BrickRed},xleftmargin=0em}
\DefineVerbatimEnvironment{Scode}{Verbatim}{fontsize=\small}
\fvset{listparameters={\setlength{\topsep}{-0.1ex}}}
\renewenvironment{Schunk}%
{\renewcommand{\baselinestretch}{0.85} \vspace{\topsep}}%
{\renewcommand{\baselinestretch}{1.00} \vspace{\topsep}}
% \renewenvironment{knitrout}
% {\renewcommand{\baselinestretch}{0.85}}
% {\renewcommand{\baselinestretch}{1.00}}

%----------------------------------------------------------------------
% The usual usefuls
% This is a file of useful extra commands snatched from
% Michael Hills, David Clayton, Bendix Carstensen & Esa Laara.
%

% Commands to draw observation lines on follow-up diagrams
%
% Horizontal lines
%

% exit time with failure, bullet
\newcommand{\hfail}[1]{\begin{picture}(250,5)
       \put(0,0){\line(0,1){2.5}}
      \put(0,0){\line(0,-1){2.5}}
      \put(0,0){\line(1,0){#1}}
      \put(#1,0){\circle*{5}}
   \end{picture}}

% exit time with censoring, open circle
\newcommand{\hcens}[1]{\begin{picture}(250,5)
         \put(0,0){\line(0,1){2.5}}
      \put(0,0){\line(0,-1){2.5}}
      \put(0,0){\line(1,0){#1}}
%      \put(#1,0){\line(0,1){2.5}}
%      \put(#1,0){\line(0,-1){2.5}}
% BxC Changed this to an open circle instead of a line
      \put(#1,0){\circle{5}}
   \end{picture}}

%
% Diagonals for Lexis diagrams
%
\newcommand{\dfail}[1]{\begin{picture}(250,250)
      \put(0,0){\line(1,1){#1}}
      \put(#1,#1){\circle*{5}}
   \end{picture}}

\newcommand{\dcens}[1]{\begin{picture}(250,250)
      \put(0,0){\line(1,1){#1}}
%      \put(#1,#1){\line(0,1){2.5}}
%      \put(#1,#1){\line(0,-1){2.5}}
% BxC Changed this to an open circle instead of a line
      \put(#1,#1){\circle{5}}
   \end{picture}}

%
% Horizontal range diagrams
%
\newcommand{\hrange}[1]{\begin{picture}(200,5)
     \put(0,0){\circle*{5}}
     \put(0,0){\line(1,0){#1}}
     \put(0,0){\line(-1,0){#1}}
   \end{picture}}

%
% Tree drawing
%
\newcommand{\tree}[3]{\setlength{\unitlength}{#1}\begin{picture}(0,0)
   \put(0,0){\line(3,2){1}}
   \put(0,0){\line(3,-2){1}}
   \put(0.81,0.54){\makebox(0,0)[br]{\footnotesize #2\ }}
   \put(0.81,-0.54){\makebox(0,0)[tr]{\footnotesize #3\ }}
\end{picture}}

\newcommand{\wtree}[3]{\setlength{\unitlength}{#1}\begin{picture}(0,0)
   \put(0,0){\line(1,1){1}}
   \put(0,0){\line(1,-1){1}}
   \put(0.8,0.8){\makebox(0,0)[br]{\footnotesize #2\ }}
   \put(0.8,-0.8){\makebox(0,0)[tr]{\footnotesize #3\ }}
\end{picture}}

\newcommand{\ntree}[3]{\setlength{\unitlength}{#1}\begin{picture}(0,0)
   \put(0,0){\line(2,1){1}}
   \put(0,0){\line(2,-1){1}}
   \put(0.8,0.4){\makebox(0,0)[br]{\footnotesize #2\ }}
   \put(0.8,-0.4){\makebox(0,0)[tr]{\footnotesize #3\ }}
\end{picture}}

%
% Other commands
%
\newcommand{\T}{\scriptsize\text T}
\newcommand{\prob}[0]{\text{\rm Pr}}
\newcommand{\nhy}[0]{_{\oslash}}
\newcommand{\true}[0]{_{\text{\rm \tiny T}}}
\newcommand{\hyp}[0]{_{\text{\rm \tiny H}}}
% \newcommand{\mpydiv}[0]{\stackrel{\textstyle \times}{\div}}
% Changed to slightly smaller symbols
\newcommand{\mpydiv}[0]{\stackrel{\times}{\scriptstyle \div}}
\newcommand{\mie}[1]{{\it #1}}
\newcommand{\mycircle}[0]{\circle*{5}}
\newcommand{\smcircle}[0]{\circle*{1}}
\newcommand{\corner}[0]{_{\text{\rm \tiny C}}}
\newcommand{\ind}[0]{\hspace{10pt}}
\newcommand{\gap}[0]{\\[5pt]}
\renewcommand{\S}[0]{section~}
\newcommand{\blank}[0]{$\;\,$}
\newcommand{\vone}{\vspace{1cm}}
\newcommand{\ljust}[1]{\multicolumn{1}{l}{#1}}
\newcommand{\cjust}[1]{\multicolumn{1}{c}{#1}}
\newcommand{\mean}{\text{\rm Mean}}
\newcommand{\transpose}{^{\mbox{\tiny T}}}
\newcommand{\histog}[5]{\rule{1mm}{#1mm}\,\rule{1mm}{#2mm}\,\rule{1mm}{#3mm}\,\rule{1mm}{#4mm}\,\rule{1mm}{#5mm}}
\newcommand{\pmiss}{P_{\mbox{\tiny miss}}}
%
% Below is BxCs commands inserted
%
\newcommand{\bc}{\begin{center}}
\newcommand{\ec}{\end{center}}

\newcommand{\bd}{\setlength{\parskip}{1ex} \begin{description}}
\newcommand{\ed}{\end{description} \setlength{\parskip}{2ex}}
\newcommand{\bdx}{\begin{description}} % Bendix' description macros
\newcommand{\edx}{\end{description}}

\newcommand{\bix}{\begin{itemize}}  % these are Bendix' itemizing macros
\newcommand{\eix}{\end{itemize}}
\newcommand{\bi}{\setlength{\parskip}{1ex} \begin{itemize}} % Esa's item macros 
\newcommand{\ei}{\end{itemize} \setlength{\parskip}{2ex}} 

\newcommand{\bn}{\begin{equation}}
\newcommand{\en}{\end{equation}}
\newcommand{\be}{\begin{enumerate}}
\newcommand{\ee}{\end{enumerate}}
\newcommand{\bes}{\begin{eqnarray*}}
\newcommand{\ees}{\end{eqnarray*}}
\newcommand{\p}{\text{\rm P}}
\newcommand{\pmat}[1]{\text{\rm P}\left\{#1\right\}}
\newcommand{\ptxt}[1]{\text{\rm P}\left\{\text{\rm #1}\right\}}
\newcommand{\E}{\text{\rm E}}
\newcommand{\V}{\text{\rm V}}
\newcommand{\BLUP}{\text{\rm BLUP}}

% \newcommand{\var}{\mbox{Var}} changed by Esa to
\newcommand{\var}{\mbox{var}}
% \newcommand{\cov}{\mbox{Cov}} changed by Esa to
\newcommand{\cov}{\mbox{cov}}
% \newcommand{\corr}{\mbox{Corr}} changed by Esa to
\newcommand{\corr}{\mbox{corr}} 

%\newcommand{\var}{\text{\rm var}}
%\newcommand{\cov}{\text{\rm cov}}
%\newcommand{\corr}{\text{\rm corr}}
\newcommand{\se}{\text{\rm s.e.}}
\newcommand{\sd}{\text{\rm std}}
\newcommand{\erf}{\text{\rm erf}}
\newcommand{\odds}{\text{\rm odds}}
\newcommand{\bin}{\text{\rm binom}}
\newcommand{\half}[1]{\frac{1}{#1}}
% \newcommand{\td}[0]{\stackrel{\textstyle \times}{\div}}
% Changed to slightly smaller symbols
\newcommand{\td}[0]{\stackrel{\scriptstyle \times}{\scriptstyle \div}}
\newcommand{\logit}{\text{\rm logit}}
\newcommand{\link}{\text{\rm link}}
\newcommand{\spn}{\text{\rm span}}
\newcommand{\OR}{\text{\rm OR}}
\newcommand{\RR}{\text{\rm RR}}
\newcommand{\ER}{\text{\rm ER}}
\newcommand{\RD}{\text{\rm RD}}
\newcommand{\AC}{\text{\rm AC}}
\newcommand{\AF}{\text{\rm AF}}
\newcommand{\PAF}{\text{\rm PAF}}
\newcommand{\SR}{\text{\rm SR}}
\newcommand{\SMR}{\text{\rm SMR}}
\newcommand{\CMF}{\text{\rm CMF}}
\newcommand{\pvp}{\text{\rm PV}$+$}
\newcommand{\pvn}{\text{\rm PV}$-$}
\newcommand{\R}{\textsf{R}}
%\newcommand{\gap}[0]{\\[5pt]} 
%\newcommand{\blank}[0]{$\;\,$}
% Conditional independence sign from Philip Dawid
\newcommand{\cip}{\mbox{$\perp\!\!\!\perp$}}

%%% Commands to comment out parts of the text
\newcommand{\GLEM}[1]{}
\newcommand{\FORGETIT}[1]{}
\newcommand{\OMIT}[1]{}

%%% Insert output from program in small text 
%%% (requires package verbatim)

\newcommand{\insout}[1]{
\scriptsize
\renewcommand{\baselinestretch}{0.8}
\verbatiminput{#1}
\renewcommand{\baselinestretch}{1.0}
\normalsize
}

% From Esa:        
%\newcommand{\T}{\text{\rm \small{T}}}
\newcommand{\id}{\text{\rm id}}
\newcommand{\Dev}{\text{\rm Dev}}
\newcommand{\Bin}{\text{\rm Bin}}
\newcommand{\probit}{\text{\rm probit}}
\newcommand{\cloglog}{\text{\rm cloglog}}
\newcommand{\EF}{\text{\rm EF}}
\newcommand{\SE}{\text{\rm SE}}
\newcommand{\IP}{\text{\rm IP}}




%----------------------------------------------------------------------
% Set up layout of pages
\oddsidemargin 1mm
\evensidemargin 1mm
\topmargin -10mm
\headheight 8mm
\headsep 5mm
\textheight 240mm
\textwidth 165mm
%\footheight 5mm
\footskip 15mm
\renewcommand{\topfraction}{0.9}
\renewcommand{\bottomfraction}{0.9}
\renewcommand{\textfraction}{0.1}
\renewcommand{\floatpagefraction}{0.9}
\renewcommand{\headrulewidth}{0.1pt}
\setcounter{secnumdepth}{4}
% \setcounter{tocdepth}{2}

%----------------------------------------------------------------------
% How to insert a figure in a .rnw file
\newcommand{\rwpre}{./graph/gr}
\newcommand{\insfig}[3]{
\begin{figure}[h]
  \centering
  \includegraphics[width=#2\textwidth]{\rwpre-#1}
  \caption{#3}
  \label{fig:#1}
% \afterpage{\clearpage}
\end{figure}}
\newcommand{\linput}[1]{
% \clearpage 
\afterpage{\hfill \ldots input from \texttt{#1.tex}} 
\fancyfoot[OR,EL]{\footnotesize \texttt{#1.tex}} 
\input{#1}}

%----------------------------------------------------------------------
% Here is the document starting with the titlepage
\begin{document}

%----------------------------------------------------------------------
% The title page
\setcounter{page}{1}
\pagenumbering{roman}
\pagestyle{plain}
\thispagestyle{empty}
% \vspace*{0.05\textheight}
\flushright
% The blank below here is necessary in order not to muck up the
% linespacing in title if it has more than 2 lines
{\Huge \bfseries \Title

}\ \\[-1.5ex]
\noindent\textcolor{blaa}{\rule[-1ex]{\textwidth}{5pt}}\\[2.5ex]
\large
\Where \\
\Dates \\
\Homepage \\
%\Version \\[1em]
\normalsize
Compiled \today,\ \currenttime\\
from: \texttt{\currfileabspath}\\[1em]
% \input{wordcount}
\normalsize
\vfill
\Faculty
% End of titlepage
% \newpage

%----------------------------------------------------------------------
% Table of contents
% \listoftables
% \listoffigures
\tableofcontents
\clearpage

%----------------------------------------------------------------------
% General text layout
\raggedright
\parindent 1em
\parskip 0ex
\cleardoublepage

%----------------------------------------------------------------------
% General page style
\pagenumbering{arabic}
\setcounter{page}{1}
\pagestyle{fancy}
\renewcommand{\chaptermark}[1]{\markboth{\textsl{#1}}{}}
\renewcommand{\sectionmark}[1]{\markright{\thesection\ \textsl{#1}}{}}
\fancyhead[EL]{\bf \thepage \quad \rm \rightmark}
\fancyhead[ER]{\rm \Tit}
\fancyhead[OL]{\rm \leftmark}
\fancyhead[OR]{\rm \rightmark \quad \bf \thepage}
\fancyfoot{}

\renewcommand{\rwpre}{./prefix}
% command to include an exercise/solution with an initial pagethrow etc.
\newcommand{\pput}[2]{\clearpage
                   \fancyhead[ER]{\sl SPE: \ExcSol}
                   \fancyhead[OL]{\sl IARC, 2018}
                   \fancyhead[OR]{\sl \rightmark \quad \bf \thepage}
                   \fancyhead[EL]{\bf \thepage \quad \sl \rightmark}
                   % \texttt{#1} ## write the name of the input file
                   \input{#1}}

\newpage
\addcontentsline{toc}{section}{Program}

\section*{Program}
\noindent
\begin{tabular}{r@{ -- }rp{13cm}}
\multicolumn{3}{l}{\bf Daily timetable} \\
 9:00 &  9:30 & Recap of yesterday's practicals \\
 9:30 & 10:30 & Lecture \\
10:30 & 11:00 & Coffee break \\
11:00 & 13:00 & Practical \\
13:00 & 14:00 & Lunch \\
14:00 & 14:30 & Recap of morning's practical \\
14:30 & 15:30 & Lecture \\
15:30 & 16:00 & Tea break \\
16:00 & 18:00 & Practical \\[2em]
\end{tabular}

\noindent
\begin{tabular}{r@{ -- }rp{13cm}}
\multicolumn{3}{l}{\bf Friday 2 June -- Day 1} \\
 8:30 &  9:00 & Registration \\
 9:00 &  9:15 & Welcome (KF) \\
 9:15 & 10:15 & R History and Ecology\\% (MP, 15 min shorter than before, remote mode) \\ 
10:15 & 10:30 & Data manipulation with \texttt{tidyverse} (DG) \\
10:30 & 11:00 & Coffeee break \\
11:00 & 13:00 & Practical---choose one: \newline
                1: Practice with basic \R\ \&  Reading data into \R\ \newline
                2: Data manipulation with \texttt{tidyverse} \\

% Students can choose:  \newline EITHER -- Practice with basic \R\ \& Simple data input (MP, 
% based on old exercises 1.1-1.2 to be trimmed according to the postmortem of 2019) for new users of \R\, \newline 
% OR -- Data manipulation with \texttt{dplyr} (DG, based on old 1.4) for \newline experienced \R\ users
% (to avoid these people getting bored in the beginning). \\
  
13:00 & 14:00 & Lunch \\
14:00 & 14:30 & Recap of morning practical \\
14:30 & 15:30 & Practical: Tabulation \& Graphics in \R\ (KF)\\ %(KF, based on old 1.3, 1.5 \newline -- especially the latter part sufficiently shortened from last time)  \\
15:30 & 16:00 & Tea break\\
16:00 & 16:30 & Simple Poisson and binary regression (JP)\\ %(JP, 30 min shorter than previously 
% $\Rightarrow$ max 15 slides! Leave at least case-controlling out?) \\
<<<<<<< HEAD
16:30 & 18:00 & Practical: Incidence rates and proportions and their contrasts by \texttt{glm} \\ 
=======
16:30 & 18:00 & Practical: Analysis of hazard rates, their ratios and differences and binary regression\\
>>>>>>> 5791fd7f29f5ff10cf0a981642ba6fefd796a94c
                % (JP, old 1.7 to be shortened but add a couple of simple items on binary regression
                % of \texttt{lowbw} on \texttt{hyp} in \texttt{births} data?) \\ 
18:00 & 21:00 & Welcome reception at Delta terrace ($4^\text{th}$ floor) \\[1em]
\end{tabular}

\noindent
\begin{tabular}{r@{ -- }rp{13cm}}
\multicolumn{3}{l}{\bf Saturday 3 June -- Day 2} \\
 9:00 &  9:30 & Recap of yesterday's practicals. \\
 9:30 & 10:00 & Introduction to splines (MP)\\ % (MP, 15 min shorter than previously, remote mode) \\
10:00 & 10:45 & Linear and generalized linear models (EL)\\ % EL, 15 min longer than before; \newline add a few things on GLMs) \\
10:45 & 11:15 & Coffeee break \\
<<<<<<< HEAD
11:15 & 13:00 & Practical: Linear, logistic \& Poisson regression  
                with estimation and reporting of linear and curved
                effects\\ %(EL, old exercises 1.9-1.10 to be revised) \\
=======
11:15 & 13:00 & Practical:  Estimation of effects: simple and more complex \& Poisson regression \& analysis of curved effects\\ %(EL, old exercises 1.9-1.10 to be revised) \\
>>>>>>> 5791fd7f29f5ff10cf0a981642ba6fefd796a94c
13:00 & 14:00 & Lunch \\
14:00 & 14:30 & Recap of morning practical \\
14:30 & 15:30 & Causal inference 1: basic concepts\\ % (KF, to be revised) \\
15:30 & 16:00 & Tea break\\
16:00 & 18:00 & Practical: Causal inference\\[1em] % (KF, old 1.14 to be revised) \\        
\end{tabular}

\noindent
\begin{tabular}{r@{ -- }rp{13cm}}
 \multicolumn{3}{l}{\bf Sunday 4 June -- Day 3} \\
 9:00 &  9:30 & Recap of yesterday's practicals \\
 9:30 & 10:30 & More advanced graphics in R, including
                \texttt{ggplot2} (MP)\\ % (MP, remote mode \newline -- or DG on the spot?)\\
10:30 & 11:00 & Coffee break. \\
11:00 & 13:00 & Practical: Graphics meccano \\% (MP, old 1.11) \\
15:00 & 17:00 & Old Town Adventure \newline \url{https://360.ee/en/team-building-games/old-town-adventure-3/}\\[1em]
\end{tabular}

\noindent
\begin{tabular}{r@{ -- }rp{13cm}}
\multicolumn{3}{l}{\bf Monday 5 June -- Day 4} \\
 9:00 &  9:30 & Recap of yesterday's practicals \\
 9:30 & 10:30 & Survival analysis: Kaplan Meier \& basic
                Cox-model, and basic analysis of competing risks (JP)\\ % (JP, perhaps some revision?)\\
10:30 & 11:00 & Coffee break\\
11:00 & 13:00 & Practical: Survival analysis with competing risks: Oral cancer patients\\ % (JP, old 1.12 to be somewhat revised)\\
13:00 & 14:00 & Lunch \\
14:00 & 14:30 & Recap of morning practical \\
14:30 & 15:30 & Dates in R; follow-up representation in \texttt{Lexis} objects,
                time-splitting,  and SMR (BxC)\\
15:30 & 16:00 & Tea break. \\
16:00 & 18:00 & Practical: Time-splitting, time-scales and SMR\\[1em] % (old exercise 1.13 to be revised, BxC)\\
% 18:30 &       & Tour of the Estonian Genome Centre(?)\\[1em]
\end{tabular}

\noindent
\begin{tabular}{r@{ -- }rp{13cm}}
\multicolumn{3}{l}{\bf Tuesday 6 June -- Day 5} \\
 9:00 &  9:30 & Recap of yesterday's practicals \\
 9:30 & 10:30 & Nested and matched cc-studies \& Case-cohort studies (KF)\\
  %(KF -- EL gave this lecture before) \\
10:30 & 11:00 & Coffee break. \\
11:00 & 13:00 & Practical: Nested case-control study and case-cohort study: \\ % (KF) old exercise 1.15 prepared by EL perhaps to be slightly revised)\\
13:00 & 14:00 & Lunch \\
14:00 & 14:30 & Recap of morning practical \\
14:30 & 15:30 & Causal inference 2: model-based estimation of causal
                contrasts (EL)\\ % \newline -- a totally new lecture)\\
15:30 & 16:00 & Tea break. \\
16:00 & 18:00 & Practical: Causal inference 2: Model-based estimation of causal estimands \\ %\newline -- a totally new exercise)\\
19:00 & 22:00 & Course dinner, Vilde \& Vine\\[1em]
\end{tabular}

\noindent
\begin{tabular}{r@{ -- }rp{13cm}}
\multicolumn{3}{l}{\bf Wednesday 7 June -- Day 6} \\
 9:00 &  9:30 & Recap of yesterday's practicals \\
 9:30 & 10:30 & Multistate models, Poisson models for rates and
                simulation of \texttt{Lexis} objects (BxC)\\
10:30 & 11:00 & Coffee break. \\
11:00 & 12:15 & Practical: Time-dependent variables and multiple states\\ % (BxC, old 1.16 to be revised)\\
12:15 & 12:45 & Recap of morning practical \\
12:45 & 13:00 & Wrap-up and farewell.\\
13:00 & 14:00 & Lunch \\
% \multicolumn{2}{l}{Afternoon} & Post-mortem (Faculty only). \\
% \multicolumn{2}{l}{Evening}   & Faculty dinner.\\
\end{tabular}
\vfill
\noindent
Further material will appear at this year's course website:\\ 
\url{http://bendixcarstensen.com/SPE}


\cleardoublepage
\pagenumbering{arabic}
\pagestyle{fancy}
\renewcommand{\sectionmark}[1]{\markboth{\thesection #1}{\thesection \ #1}}
\renewcommand{\headrulewidth}{0.1pt}
\setcounter{section}{0}
\setcounter{page}{1}

%% \end{document}

\chapter{Exercises}
\newcommand{\ExcSol}{Exercises}

\addcontentsline{toc}{section}{Introduction to practicals}
% \section*{Introduction to practicals}

\subsection*{Data sets}

Datasets for the practicals in this course, as well as some useful R
scripts, will be available on the course homepage, in
\url{http://BendixCarstensen.com/SPE/data}. This is where you will
also find the ``housekeeping'' scripts designed to save you typing.

To get all files in one go, download the zip
file \url{http://spe-r.github.io/data.zip}.

\subsection*{Graphical User Interfaces to R}

When running the exercises it is a good idea to use a text editor
instead of typing your commands directly at the R prompt. On Windows
and macOS, R comes with a basic graphical user interface including a
built-in text editor. Many people like to use the RStudio interface to
R, which includes a very powerful syntax-highlighting editor.

\subsection*{Keyboard shortcuts}

In the past we have found that some participants have had difficulty
finding keys for symbols that are commonly used in the R language.  In
particular, the tilde symbol \verb+~+ is used in all modelling
functions but not directly available on some keyboards. If this
affects you then please consult the Wikipedia
page: \url{http://en.wikipedia.org/wiki/Tilde#Keyboards} for advice on
the combination of key presses you will need to get tilde.

\subsection*{Recaps}
The R-scripts used during the course for the recaps will be available
in \url{http://BendixCarstensen.com/SPE/recap}.

\subsection*{Ask for help}

The faculty are here to help you. Ask them for help.




% Thursday 1
\pput{basic-e}{Practice with basic \R}
\pput{dinput-e}{Simple reading and data input}
\pput{tab-e}{Tabulation}
\pput{dplyr-e}{Data manipulation with \textt{dplyr}}
\pput{graph-intro-e}{Introduction to graphs in \R}
\pput{simulation-e}{Simple simulation}

% Friday 2
\pput{rates-rrrd-e}{Rates, rate ratio and rate difference with \texttt{glm}}  % BxC / EL
\pput{logistic-e}{Logistic regression with \texttt{glm}}  % EL
\pput{effects-e}{Simple estimation of effects}
\pput{cont-eff-e}{Estimation and reporting of linear and curved effects}  % BxC

% Saturday 3
\pput{graphics-e}{Graphical meccano}                 % MP

% Sunday 4
\pput{oral-e}{Survival and competing risks in oral cancer}  % EL
\pput{DMDK-e}{Time-splitting and SMR (Danish diabetes patients)} % BxC

% Monday 5
\pput{causal-e}{Simulation and causal inference}     % KF
\pput{occoh-caco-e}{Nested case-control and case-cohort study} % EL

% Tuesday 6
\pput{renal-e}{Multistate-model: Renal complications} % BxC

%% check if withoutsolutions variable has been pass when compiling 
%% and stop or continue the compilation accordingly
\ifdefined\withsolutions
  
\else
  \end{document}
\fi

% \end{document}

\chapter{Solutions}
\renewcommand{\ExcSol}{Solutions}

There is a chapter for each of the exercises used at the
course. This is either a printout of the R-program that
performs the analyses, as well as the graphs produced by the
programs, or output from an R-weave solution file with a bit more
elaborate text.

The code and the output from these programs are also available from
the course homepage in \url{http://BendixCarstensen/SPE/R}; they are
called \texttt{xxx-s.R}; just before each chapter you will find a line
with the text \texttt{xxx-s}, indicating that the name of the script
will be \texttt{xxx-s.R}.

% Note that exercsises without sultion are replaced by an \addtocounter
% in order to syncronize the numbering between the exercise chapter
% and the solution chapter

% Thursday 1
\addtocounter{section}{1}
% \pput{basic-s}{Practice with basic \R}
\addtocounter{section}{1}
% \pput{dinput-s}{Simple reading and data input}
\pput{tab-s}{Tabulation}
\pput{dplyr-s}{Data manipulation with \textt{dplyr}}
\pput{graph-intro-s}{Introduction to graphs in \R}  % KF
\pput{simulation-s}{Simple simulation}  % KF
% 
% Friday 2
\pput{rates-rrrd-s}{Rates, rate ratio and rate difference with \texttt{glm}}  % JP
\pput{logistic-s}{Logistic regression with \texttt{glm}}  % JP
\pput{effects-s}{Simple estimation of effects}  % EL
\pput{cont-eff-s}{Estimation and reporting of linear and curved effects}  % EL

% Saturday 3
\addtocounter{section}{1}
% \pput{graphics-s}{Graphical meccano} % MP

% Sunday 4
\pput{oral-s}{Survival and competing risks in oral cancer}  % EL
\pput{DMDK-s}{Time-splitting and SMR (Danish diabetes patients)} % BxC

% Monday 5
\pput{causal-s}{Simulation and causal inference}     % KF
\pput{occoh-caco-s}{CC study: Risk factors for Coronary heart disease} % EL

% Tuesday 6
\pput{renal-s}{Multistate-model: Renal complications} % BxC

\end{document}
